% C.42 candidate for integration after root's final proof read.
\begin{statement}{C.42}{Proved}{The three-primary part of the iterated finite lifting obstruction vanishes}
Let $G$ be the actual finite group of order $2^{146}3^{30}$ in the
retained iterated-descent calculation, with its original $176$-generator
pc presentation.  The complete presentation, the two modular generators
$s,r$, their prescribed images $a,b$, and the left peripheral frames
$B_s,B_r$ are specified in
\texttt{pbc\_iterated\_shortcut\_v2\_original\_source.json}.
Write $G_n$ for the actual quotient obtained by killing pc generators
with index greater than $n$, for the normal cuts used below.
Then the actual quotient map induces an isomorphism
\[
 H_2(G,\mathbb Z_3)\xrightarrow{\ \sim\ }H_2(G_{26},\mathbb Z_3),
 \qquad 9H_2(G,\mathbb Z_3)=0.
 \tag{C.42.1}
Together with the previously retained two-primary bound, this gives the
integral annihilator
\[
                 2^{14}3^2H_2(G,\mathbb Z)=0.
 \tag{C.42.2}

Let $L$ be the free group on the original pc generators, let
$R=\ker(L\to G)$, and put $K=R/[L,R]$.  The central framed-product
defect $T\in K$ is the sum of the three central differences
\[
 (B_ss)-(aB_s),\qquad (B_rr)-(bB_r),\qquad (ab)-(sr).
 \tag{C.42.3}
Its three-primary component in the Hopf subgroup is zero:
\[
                T_{(3)}=0\quad\hbox{in }H_2(G,\mathbb Z)_{(3)}.
 \tag{C.42.4}
No boundary correction is needed for this component.
This is a proved component of the iterated finite lifting problem.
It does not assert the vanishing of the two-primary defect or the
existence of a full gentle lift realizing a third independent
iterated-parity value.
\end{statement}
\begin{proof}
The source is identified throughout by its SHA256 hash
\texttt{95fb384593bfaf97d468bbc51f11cd797484a08cb3ec3b954961085b520ffc88},
retained with the executed certificates.  Its provenance is the actual
finite word image, rather than a group selected only by its order or
composition factors.  The new checker reads all $176$ original power
relations and all $15\,400$ original conjugation relations.  It verifies
the normal suffixes and obtains the following actual quotient maps:
\[
\begin{array}{c|c}
 \text{map}&\text{kernel}\\\hline
 G_{176}\to G_{104}&(C_2)^{72}\\
 G_{104}\to G_{86}&(C_3)^{18}\\
 G_{86}\to G_{26}&\text{a }2\text{-group of order }2^{60}\\
 G_{26}\to G_{17}&(C_3)^9\\
 G_{17}\to G_4&\text{a }2\text{-group of order }2^{13}.
\end{array}
 \tag{C.42.5}
The two unspecified $2$-groups need not be abelian.  Each elementary
abelian assertion is checked from every power and internal conjugation
relation in its actual quotient.  Quotienting by a finite normal
$2$-group induces an isomorphism on all homology with $\mathbb Z_3$
coefficients, by the homology Hochschild--Serre spectral sequence.

\paragraph{The actual module calculations.}
For $A_d=\ker(G_h\to G_l)$ with $(l,h,d)=(17,26,9)$ or
$(86,104,18)$, take the original pc generators with indices
$l+1,\ldots,h$ as the $\mathbb F_3$-basis.  The action matrices are
read directly from their original conjugation words.  For $d=9$, the
normal $2^{13}$ subgroup of $G_{17}$ has coinvariants zero on $A_9$;
for $d=18$, the normal $2^{60}$ subgroup of $G_{86}$ has coinvariants
zero on $A_{18}$.  In both cases this is a full-rank calculation of
the stacked original matrices $M_g-I$, with ranks $9$ and $18$.
Since the normal subgroup has order prime to three, averaging and
Hochschild--Serre give
\[
                   H_i(G_l,A_d)=0\qquad(i\ge0).
 \tag{C.42.6}

For the exterior-square module, the action rows are the actual
two-by-two minors of those same matrices.  The complete rank data are
\[
\begin{array}{c|r|r|r|r}
 d&\dim\Lambda^2A_d&\text{rows}&\text{rank}
                   &\dim(\Lambda^2A_d)_{G_l}\\\hline
 9&36&612&35&1\\
 18&153&13\,158&153&0.
\end{array}
 \tag{C.42.7}
Every row is checked.  The certificate additionally retains the selected
independent original rows and verifies their ranks with a second
plain-Python elimination.  In particular, the zero in the last entry
is an exact finite certificate, not an inferred representation type.

For an elementary abelian three-kernel $A$ with quotient $Q$, the
total-degree-two terms of the homology spectral sequence are
\[
 H_2(Q,\mathbb Z_3),\qquad H_1(Q,A),\qquad
                         H_0(Q,\Lambda^2A).
 \tag{C.42.8}
Here $H_2(A,\mathbb Z_3)=\Lambda^2A$ follows from the integral
abelian-group identity.  For $A_{18}$ the last two terms vanish.
The only possible outgoing differential from the first term is
$d_2$ to $A_{18,Q}$, which vanishes by (C.42.6); there is no incoming
differential and no later outgoing differential.  Thus the actual map
$G_{104}\to G_{86}$ induces an isomorphism on $H_2(-,\mathbb Z_3)$.
Combining this with the two normal $2$-kernel maps proves the
isomorphism in (C.42.1).

For $A_9$, the middle term in (C.42.8) vanishes and the last is killed
by three.  Its spectral-sequence filtration therefore multiplies a
bound for $H_2(G_{17},\mathbb Z_3)$ by at most three.  No splitting
of that filtration is assumed.

\paragraph{An elementary bound for the actual order-$54$ quotient.}
The four-generator quotient has relative orders $(3,2,3,3)$ and
trivial power words.  Its actual subgroup $A=\langle g_3,g_4\rangle$
is $C_3^2$, with quotient $C_6$.  On $A$, the actions of $g_1,g_2$
are respectively
\[
             U=\begin{pmatrix}1&0\\1&1\end{pmatrix},
             \qquad -I.
 \tag{C.42.9}
These matrices and the quotient presentation are checked directly in
the original relations.  A generator of $C_6$ acts by $-U$, so its
action minus the identity is invertible over $\mathbb F_3$.
The cyclic resolution gives $H_1(C_6,A)=0$, while
$H_2(C_6,\mathbb Z_3)=0$.  Consequently $H_2(G_4,\mathbb Z_3)$ is
a subquotient of $(\Lambda^2A)_{C_6}$ and is killed by three.
The normal $2^{13}$ comparison and the $A_9$ filtration now give
$9H_2(G_{26},\mathbb Z_3)=0$.  This proves (C.42.1).
For completeness, the retained actual derived-series calculation has
successive abelian-factor exponents
\[
                         6,\ 3,\ 4,\ 6,\ 4,\ 3,\ 2.
\]
For an abelian normal kernel of exponent $r$, the last two terms of
(C.42.8), with that kernel in place of $A$, are killed by $r$.
The total-degree-two filtration therefore multiplies an annihilator
of the quotient's multiplier by at most $r^2$.  Induction through
the derived series gives the annihilator
$(6\cdot3\cdot4\cdot6\cdot4\cdot3\cdot2)^2=2^{14}3^8$.
The complete derived generators and their actual abelian invariants
are retained in \texttt{pbc\_derived\_h2\_bound.json}.
In particular the two-primary bound is $2^{14}$.  Combining it
with the sharper three-primary calculation, and using that finite-group
positive homology has no other primary parts here, proves (C.42.2).

An independently executed GAP calculation of the actual order-$54$
presentation returns its full multiplier $C_3$.  The elementary
argument above suffices for the bound, so it does not depend on
that software result.  The relevant documented multiplier and
source-centralization conventions are
\cite[\S7.9.2,\S7.9.4]{PBC_Polycyclic2026}.

\paragraph{A literal original-row identity for the defect.}
The quotient has $26$ pc generators and $351$ formal central tails.
All three terms in (C.42.3), and all original finite centralizer
words, are recollected in this actual quotient presentation.
Raw tail vectors from the $176$-generator collector are not
truncated: only the group words are projected before recollection.
In the resulting original tail coordinates the target is
\[
                             T_{26}=e_{12}+e_{24}.
 \tag{C.42.10}
The exact search modulo nine uses only the nine genuine
$r$-centralizer commutators.  Its $16\,278$-node arithmetic DAG
is independently expanded into original leaves.  Every reduced
coordinate bound $g=\gcd(d,9)$ is justified by explicit integers
$a,b$ with $ad+9b=g$, where $d e_i$ is itself a retained integer
combination of original power-consistency identities.

The final displayed certificate removes every coordinate and
modulus leaf before forming its integer sum.  It then applies the
explicit original power identities just described.  The exported
ledger contains exactly $440$ original identity rows and one
genuine commutator row, and verifies the literal equality
\[
       T_{26}+3c_{74}+\sum_{j=1}^{440} n_j R_j
                              =9z\quad\text{in }\mathbb Z^{351}.
 \tag{C.42.11}
Every $n_j$, every original row $R_j$, and all $351$ coordinates
of $z$ are retained in
\texttt{pbc\_quotient26\_literal\_original\_identity.json}.
There are no coordinate or modulus rows in this final witness.
The standard GAP collector independently replays all $442$ rows
appearing in it, including the target and the commutator, in the
original unprocessed presentation.  A separate earlier replay
checks the full set of $443$ used leaves before consolidation.
The source hashes and both complete successful outputs are retained.

The index $74$ is the first of the nine original $r$-centralizer
corrections; the two power-word positions and every $s$-correction
have coefficient zero.  In the actual source centralization,
$c_{74}$ is central and has exponent dividing three, because it
is the commutator of a lift of an element centralizing $r$ with a
lift of $r$, and $r^3$ is central.  Thus $3c_{74}=0$ integrally.
Equation (C.42.11) proves $T_{26}\in9K_{26}$.

Finally the image of (C.42.3) in the free abelianization is zero:
the three differences telescope.  The Hopf exact sequence therefore
places $T$ and $T_{26}$ in the respective finite torsion subgroups
$H_2(G,\mathbb Z)$ and $H_2(G_{26},\mathbb Z)$ of the centralization
kernels.  Since the three-primary subgroup is killed by nine, its
map to $K_{26}/9K_{26}$ is injective.  Indeed a finitely generated
abelian group splits abstractly as a free group plus its finite
torsion subgroup, and nine kills precisely the relevant component.
Hence $(T_{26})_{(3)}=0$.  Hopf naturality identifies this with the
image of $T_{(3)}$ under the actual isomorphism (C.42.1), proving
(C.42.4).  The remaining two-primary component is not addressed
by this last inference.
\end{proof}
